\chapter*{2024年\textsf{\textsc{Hdu}}「线性代数」期中试题}

\ExplSyntaxOn
\vspace{-12.0pt}\setcounter{page}{1}
\begin{center}
  \begin{minipage}{.66\textwidth}
    \centering
    \begin{tasks}
      [ label-width = 5\ccwd, item-indent = 2cm,
        after-item-skip = 1ex ](2)
      \task[\textbf{课程信息}：] 理学院/A0714040
      \task[\textbf{解析提供}：] \textsf{幻想大白}
      \task[\textbf{授课教师}：] 大学数学教学团队
      \task[{\makebox[2em][c]{\textbf{\textsf\TeX}}}\textbf{排版}：] \textsf{Izumi Sakai}
    \end{tasks}
  \end{minipage}
  \begin{minipage}{.15\textwidth}
    \centering
    \begin{tikzpicture}
      \node [ inner~sep = 0 ] (wechat) {
        \fancyqr
          [ size = 4em, image~padding = 1,
            image = \raisebox{4em/33}{\textcolor{DarkGreen \l__whats_darkmode_fading_tl}{\faIcon{weixin}}},
            r~color = DarkGreen \l__whats_darkmode_fading_tl, version = 4
          ]{http://weixin.qq.com/r/sRJiemTEGjllrSAZ90ej}
        };
      \node [ inner~sep = .5ex, align = center, right, font = \scriptsize ]
        at (wechat.east) {未\\央\\学\\社};
    \end{tikzpicture}
  \end{minipage}
  \begin{minipage}{.15\textwidth}
    \centering
    \begin{tikzpicture}
      \node [ inner~sep = 0 ] (qq) {
        \fancyqr
          [ size = 4em, image~x~padding = 2, image~y~padding = 1,
            image = \raisebox{4em/33}{\textcolor{MidnightBlue \l__whats_darkmode_fading_tl}{\faIcon{qq}}},
            l~color = MidnightBlue \l__whats_darkmode_fading_tl, version = 4
          ]{https://qm.qq.com/q/rbfYuqPhQs}
        };
      \node [ inner~sep = .5ex, align = center, left, font = \scriptsize ]
        at (qq.west) {数\\学\\营};
    \end{tikzpicture}
  \end{minipage}
\end{center}
\ExplSyntaxOff

\section{选择题（每题3分，共18分）}

\begin{problem}
  设 $\mat A = \diag{(a_1,~ a_2,~ \cdots,~ a_n)},~
      a_i \neq a_j~ (i,~ j = 1,~ 2,~ \cdots,~ n)$，
  那么与 $\mat A$ 乘法可交换的矩阵为
  \begin{tasks}(4)
    \task[\true] 对角矩阵
    \task 上三角矩阵
    \task 下三角矩阵
    \task 对称矩阵
  \end{tasks}
\end{problem}
\scratch
\begin{solution}
  $\mat A$ 为对角阵，若 $\mat{AB} = \mat{BA}$，
  显然 $\mat B$ 也为对角阵，所以A选项正确.
\end{solution}

\begin{problem}
  如果 $n$ 阶方阵 $\mat A$ 满足 $\trans{\mat A} \mat A =
  \mat A \trans{\mat A} = \mat E$，则 $\mat A$ 的行列式有
  \begin{tasks}(4)
    \task $\det A = 0$
    \task[\true] $\det A \neq 0$
    \task $\det A = k$（$k > 1$）
    \task $\det A = k$（$k < 1$）
  \end{tasks}
\end{problem}
\scratch
\begin{solution}
  由 $\mat A \trans{\mat A} = \mat E$ 得
  $\det{A \trans A} = \det E$，所以
  $\det A^2 = 1,~ \mat A = \pm1 \neq 0$.
\end{solution}

\begin{problem}
  设 $\mat A$ 为 $n$ 阶对称矩阵，$B$ 为 $n$ 阶反对称矩阵，则下列矩阵中为对称矩阵的是
  \begin{tasks}(4)
    \task[\true] $\mat{BAB}$
    \task $\mat{ABA}$
    \task $\mat{ABAB}$
    \task $\mat{BABA}$
  \end{tasks}
\end{problem}
\scratch
\begin{solution}
  由题意
  $\mat A = \trans{\mat A},~ \mat B = -\trans{\mat B}$.
  A选项
  $ \mat{BAB} = \trans{\mat B} \trans{\mat A} \trans{\mat B}
  = (-\mat B) \cdot (\mat A) \cdot (-\mat B) = \mat{BAB}$. 所以A选项正确.
\end{solution}

\begin{problem}
  设 $\mat A$ 为 $n$ 阶方阵，$\det A = a \neq 0$，
  $\mat A^*$ 为 $\mat A$ 的伴随矩阵，则 $\det{A^*} = $
  \begin{tasks}(4)
    \task $a$
    \task $a^{-1}$
    \task[\true] $a^{n-1}$
    \task $a^*$
  \end{tasks}
\end{problem}
\scratch
\begin{solution}
  $\det{A^*} = \det A^{n-1} = a^{n-1}$. 所以C选项正确.
\end{solution}

\begin{problem}
  设 $\mat A,~ \mat B$ 均为 $n$ 阶方阵，则下列说法正确的是
  \begin{tasks}(2)
    \task 若 $\mat{AB} = \mat O$，
    则 $\mat A = \mat O$ 或 $\mat B = \mat O$
    \task 若 $\mat{AB} = \mat{AC}$，且 $\mat A \neq 0$，
    则 $\mat B = \mat C$
    \task[\true] 若 $\mat{AB} = \mat O$，则 $\det A = \mat O$
    或 $\det B = \mat O$
    \task 若 $\mat{AB} = \mat{CA}$，且 $\mat A$ 可逆，
    则 $\mat B = \mat C$
  \end{tasks}
\end{problem}
\scratch
\begin{solution}
  由 $\mat AB = \mat O$ 得 $\det{AB} = \det A \det B = \det O$，
  所以 $\det A = 0$ 或 $\det B = 0$，C选项正确.
\end{solution}

\begin{problem}
  设三阶实的方阵 $\mat A = (a_{ij})$ 满足
  $\mat A^* = \trans{\mat A}$，$\mat A^*$ 为 $\mat A$的伴随矩阵，
  $\trans{\mat A}$ 为 $\mat A$ 的转置矩阵.
  若 $a_{11} = a_{12} = a_{13} > 0$，则 $a_{11} = $
  \begin{tasks}(4)
    \task $\sqrt2/2$
    \task $\sqrt3/2$
    \task $\sqrt2/3$
    \task[\true] $\sqrt3/3$
  \end{tasks}
\end{problem}
\scratch
\begin{solution}
  由 $|\mat A^*| = \det{\trans A}$ 得
  $\det A^{n-1} = \det A,~ \det A = 1$.
  将行列式按行展开得 $a_{11}^2 + a_{12}^2 + a_{13}^2 = 3a^2 = 1,~ a = 1/\sqrt3$.
\end{solution}

\section{填空题（每题3分，共12分）}

\begin{problem}
  排列 $135798642$ 是 \result{偶} （填奇或偶）排列.
\end{problem}
\begin{solution}
  逆序有 $(3,2),~ (5,2),~ (5,4),~ (7,2),~ (7,4),~ (7,6),~(9,2),~ (9,4),~
  (9,6),~ (9,8)$，共 $12$ 个，偶排列.
\end{solution}

\begin{problem}
  $ \mat E_3(1,2)
    \begin{bmatrix}
      2 & 1 & 0\\
      1 & 1 & 1\\
      3 & 0 & 2
    \end{bmatrix} = $
  \result
    {\fbox{
      $\begin{bmatrix}
      1 & 1 & 1\\
      2 & 1 & 0\\
      3 & 0 & 2
      \end{bmatrix}$}
    }.
\end{problem}

\begin{problem}
  当 $n$ 为奇数时，$n$ 阶行列式 $D =
  \begin{vmatrix}
    0       & a_{12}  & \cdots & a_{1n}\\
    -a_{12} & 0       & \cdots & a_{2n}\\
    \vdots  & \vdots  & \ddots & \vdots \\
    -a_{1n} & -a_{2n} & \cdots & a_{nn}
  \end{vmatrix} = $
  \result{$0$}.
\end{problem}
\begin{solution}
  $ D = \trans D =
    \begin{vmatrix}
      0      & -a_{12} & \cdots & -a_{1n}\\
      a_{12} & 0       & \cdots & -a_{2n}\\
      \vdots & \vdots  & \ddots & \vdots \\
      a_{1n} & a_{2n}  & \cdots & -a_{nn}
    \end{vmatrix} = (-1)^n
    \begin{vmatrix}
      0      & a_{12} & \cdots & a_{1n}\\
      a_{12} & 0      & \cdots & a_{2n}\\
      \vdots & \vdots & \ddots & \vdots \\
      a_{1n} & a_{2n} & \cdots & a_{nn}
    \end{vmatrix} = (-1)^n \trans D = -\trans D$，所以 $D = 0$.
\end{solution}

\begin{problem}
  已知三阶方阵 $\mat A$ 的第二行元素为 $1,~ 2,~ 3$，
  $\mat A$ 中第一行元素的余子式为 $3,~ a,~ 1$，则 $a = $ \result{$3$}.
\end{problem}
\begin{solution}
  $3 \times 1 + (-a) \times 2 + 1 \times 3 = 0$，所以 $a = 3$.
\end{solution}

\section{解答题A（每题6分，共24分）}

\begin{problem}
  求行列式 $D =
  \begin{vmatrix}
    a_1 & 0   & 0   & b_1\\
    0   & a_2 & b_2 & 0\\
    0   & b_3 & a_3 & 0\\
    b_4 & 0   & 0   & a_4
  \end{vmatrix}$ 的值.
\end{problem}
\scratch
\begin{solution}
  \begin{align*}
    D & = a_1
    \begin{vmatrix}
      a_2 & b_2 & 0\\
      b_3 & a_3 & 0\\
      0   & 0   & a_4
    \end{vmatrix} - b_4
    \begin{vmatrix}
      0   & a_2 & b_2\\
      0   & b_3 & a_3\\
      b_4 & 0   & 0
    \end{vmatrix}
        = a_1(a_2a_3a_4 - b_2b_3b_4) - b_1(a_2a_3b_4 - b_2b_3b_4)\\
      & = a_1a_4(a_2a_3 - b_2b_3) - b_1b_4(a_2a_3 - b_2b_3)
        = (a_1a_4 - b_1b_4)(a_2a_3 - b_2b_3)
  \end{align*}
\end{solution}

\begin{problem}
  求矩阵 $\mat A =
  \begin{bmatrix}
    2 & 5 & 0 & 0\\
    3 & 8 & 0 & 0\\
    0 & 0 & 1 & 3\\
    0 & 0 & 0 & 2
  \end{bmatrix}$ 的逆矩阵.
\end{problem}
\scratch
\begin{solution}
  令 $\mat A_1 = \begin{bmatrix} 2 & 5\\ 3 & 8\end{bmatrix},~
      \mat A_2 = \begin{bmatrix} 1 & 3\\ 0 & 2\end{bmatrix}$，
  则 $\mat A =
  \begin{bmatrix} \mat A_1 & \mat O\\ \mat O & \mat A_2 \end{bmatrix}$，
  $\mat A^{-1} =
  \begin{bmatrix}
    \mat A_1^{-1} & O\\
    \mat O        & \mat A_2^{-1}
  \end{bmatrix}$.\par\noindent
  其中 $\mat A_1^{-1} =
  \begin{bmatrix} 8 & -5\\ -3 & 2 \end{bmatrix},~
  \mat A_2^{-1} =
  \begin{bmatrix} 1 & -3/2\\ 0 & 1/2 \end{bmatrix}$. 所以 $\mat A^{-1} =
  \begin{bmatrix}
    8  & -5 & 0 & 0\\
    -3 & 2  & 0 & 0\\
    0  & 0  & 1 & -3/2\\
    0  & 0  & 0 & 1/2
  \end{bmatrix}$.
\end{solution}

\begin{problem}
  若 $\det A = \vxmat a33 = m,~
  \det B =
  \begin{vmatrix}
    a_{11} & a_{12} & b_{13}\\
    a_{21} & a_{22} & b_{23}\\
    a_{31} & a_{32} & b_{33}
  \end{vmatrix} = n$，求 $\det{A + B}$ 的值.
\end{problem}
\scratch
\begin{solution}
  将矩阵 $\mat A$、$\mat B$ 列分块，
  令 $\mat A = (\bm\alpha_1, \bm\alpha_2, \bm\alpha_3),~
  \mat B = (\bm\alpha_1, \bm\alpha_2, \bm\beta)$.
  则 $ \det{A + B} = |(2\bm\alpha_1, 2\bm\alpha_2, \bm\alpha_3 + \bm\beta)|
      = 4|(\bm\alpha_1, \bm\alpha_2, \bm\alpha_3 + \bm\beta)|
      = 4[|(\bm\alpha_1, \bm\alpha_2, \bm\alpha_3)|
        + |(\bm\alpha_1, \bm\alpha_2, \bm\beta)|] = 4(m + n)$.
\end{solution}

\begin{problem}
  已知从变变量 $x_1,~ x_2,~ x_3$ 到变量 $y_1,~ y_2,~ y_3$ 的线性变换为
  $\begin{cases}
    y_1 =  x_1 + 2x_2 - x_3\\
    y_2 = 3x_1 + 4x_2 - 2x_3\\
    y_3 = 5x_1 - 4x_2 + x_3
  \end{cases}$，求从变量 $y_1,~ y_2,~ y_3$ 到变量
  $x_1,~ x_2,~ x_3$ 的线性变换.
\end{problem}
\scratch
\begin{solution}
  令 $\mat X = \trans{(\bm x_1, \bm x_2, \bm x_3)},~
  \mat Y = \trans{(\bm y_1, \bm y_2, \bm y_3)},~
  \mat A =
  \begin{bmatrix}
    1 & 2 & -1\\
    3 & 4 & -2\\
    5 & -4& 1
  \end{bmatrix},~ \det A = 2 \neq 0$，$\mat A$ 可逆.
  由 $\mat Y = \mat{AX}$ 得
  \[
    \mat X = \mat{A}^{-1}\mat Y =
    \begin{bmatrix}
      -2    & 1 & 0\\
      -13/2 & 3 & -1/2\\
      -16   & 7 & 1
    \end{bmatrix}
    \begin{bmatrix}
      y_1\\ y_2\\ y_3
    \end{bmatrix} \longrightarrow
    \begin{cases}
      x_1 = -2y_1 + y_2\\
      x_2 = -\frac{13}{2}y_1 + 3y_2 - \frac12y_3\\
      x_3 = -16y_1 + 7y_2 + y_3
    \end{cases}
  \]
\end{solution}

\section{解答题B（每题7分，共28分）}

\begin{problem}[本题7分]
  设行列式 $D =
  \begin{vmatrix}
    2 & 1  & 4 & 1\\
    3 & -1 & 2 & 1\\
    1 & 2  & 3 & 2\\
    5 & 0  & 6 & 2
  \end{vmatrix}$，求 $M_{13} - M_{23} + 2M_{33}$，
  其中 $M_{ij}$ 是 $D$ 中第 $i$ 行 第 $j$ 列元素的余子式.
\end{problem}
\scratch[2ex]
\begin{solution}
  $ M_{13} - M_{23} + 2M_{33} = A_{13} + A_{23} + 2A_{33} + 0 \cdot A_{43} =
    \begin{vmatrix}
      2 &  1 & 1 & 1\\
      3 & -1 & 1 & 1\\
      1 & 2  & 2 & 2\\
      5 & 0  & 0 & 2
    \end{vmatrix} = -12$.
\end{solution}

\begin{problem}[本题7分]
  已知 $\mat A =
  \begin{bmatrix}
    2 & 1  & 1\\
    3 & 1  & 2\\
    1 & -1 & 0
  \end{bmatrix}$，且 $\mat{AX} + 2\mat E = \mat A^2 + \mat A + \mat X$，
  求 $\mat X$.
\end{problem}
\scratch[2ex]
\begin{solution}\leavevmode
  \begin{itemize}
    \item $\mat{AX} + 2\mat E = \mat A^2 + \mat A + \mat X$ 可改写为
    $ (\mat A - \mat E)\mat X = (\mat A - \mat E)(\mat A + 2\mat E),~
      (\mat A - \mat E)(\mat X - \mat A - 2\mat E) = \mat O$.
    \item 由于 $\det{A - E} \neq 0$，所以 $\mat A$ 可逆.
    齐次线性方程组 $\mat A\bm x = \bm 0$ 仅有零解.
    \item 所以 $\mat X - \mat A - 2\mat E = \mat O,~ \mat X = \mat A + 2\mat E =
    \begin{bmatrix}
      4 & 1 & 1\\
      3 & 3 & 2\\
      1 & -1 & 2
    \end{bmatrix}$.
  \end{itemize}
\end{solution}

\begin{problem}[本题7分]
  求 $f(x) =
  \begin{vmatrix}
    x - 2 & -2  & 2\\
    -2    & x-5 & 4\\
    2     & 4   & x-5
  \end{vmatrix} = 0$ 的根.
\end{problem}
\scratch[2ex]
\begin{solution}
  \[
      f(x) =
      \begin{vmatrix}
        x - 2 & 0     & 2\\
        -2    & x - 1 & 4\\
        4     & 0     & x - 9
      \end{vmatrix}
    = (x - 1)(-1)^{2+2}
      \begin{vmatrix}
        x - 2 & 2\\
        4     & x - 9
      \end{vmatrix}
    = (x - 1)^2(x - 10) = 0,~ x_1 = x_2 = 1,~ x_3 = 10
  \]
\end{solution}

\begin{problem}[本题7分]
  设线性方程组
  $\begin{cases}
    x_1  + 2x_2        + 3x_3 = 1\\
    2x_1 + 2x_2        + 5x_3 = 2\\
    3x_1 + \lambda x_2 + x_3  = 1
  \end{cases}$
  有唯一解，求 $\lambda$ 的取值范围。
\end{problem}
\scratch[2ex]
\begin{solution}
  令 $\mat A =
  \begin{bmatrix}
    1 & 2       & 3\\
    2 & 2       & 5\\
    3 & \lambda & 1
  \end{bmatrix},~
  \bm\beta =
  \begin{bmatrix}
    1\\ 2\\ 1
  \end{bmatrix}$. $\mat{A}\bm x = \beta$ 有唯一解，则
  $\mathbf R(\mat A) = \mathbf R(\mat A,\bm\beta) = 3,~
    \det A = \lambda + 10\neq 0,~ \lambda \neq -10$.
\end{solution}

\section{综合证明题（共18分）}

\begin{problem}[本题10分]
  设 $m$ 次多项式
  $f(x) = a_mx^m + a_{m - 1}x^{m - 1} + \cdots + a_1x + a_0$，
  $\mat A$ 为 $n$ 阶单位矩阵，定义矩阵多项式
  $f(\mat A) = a_m\mat A^m + a_{m - 1}\mat A^{m - 1} + \cdots
             + a_1\mat A + a_0\mat E$.
  若 $\bm\alpha = \trans{(1,2,3)},~ \beta = (1, \frac12, \frac13),~
  \mat A = \bm\alpha\bm\beta$，则
  \begin{tasks}[ label = \arabic*. ](2)
    \task 求 $\mat A^m$.
    \task 若 $f(x) = x^2 - x - 1$，求矩阵多项式 $f(\mat A)$.
  \end{tasks}
\end{problem}
\scratch[2ex]
\begin{solution}\leavevmode
  \begin{enumerate}
    \item
    $ \mat A^m = (\bm\alpha\bm\beta)^m
    = \bm\alpha(\bm\beta\bm\alpha)^{m-1}\bm\beta$，其中
    $\bm\alpha\bm\beta =
    \begin{bmatrix}
      1 & 1/2 & 1/3\\
      2 & 1       & 2/3\\
      3 & 3/2 & 1
    \end{bmatrix},~ \bm\beta\bm\alpha = 1 + 1 + 1 = 3$.
    所以 $\mat A^m = 3^{m - 1}
    \begin{bmatrix}
      1 & 1/2 & 1/3\\
      2 & 1   & 2/3\\
      3 & 3/2 & 1
    \end{bmatrix}$.
    \item $\mat A^2 = 3\mat A,~
      f(\mat A) = 3\mat A - \mat A - \mat E = 2\mat A - \mat E
    = \begin{bmatrix}
      1 & 2 & 2/3\\
      4 & 1 & 4/3\\
      6 & 3 & 1
    \end{bmatrix}$.
  \end{enumerate}
\end{solution}

\begin{problem}[本题8分]
  设 $n$ 阶方阵 $\mat A$ 和 $\mat B$ 满足等式 $2(\mat{A + B} = \mat{BA})$.
  \begin{tasks}[ label = \arabic*. ](2)
    \task 证明：$\mat A - 2\mat E$ 可逆.
    \task 若 $\mat B =
    \begin{bmatrix}
      1 & 1 & 0\\
      2 & 1 & 0\\
      0 & 0 & 3
    \end{bmatrix}$，求 $\mat A$.
  \end{tasks}
\end{problem}
\scratch[2ex]
\begin{solution}\leavevmode
  \begin{enumerate}
    \item \begin{proof}
      由 $2\mat A + 2\mat B = \mat{BA}$ 得
      $ 2\mat A = \mat B(\mat A - 2\mat E),~
        \mat A - 2\mat E = \frac12\mat B(\mat A - 2\mat E) - 2\mat E,~
        \longrightarrow
        \ab(\frac14\mat B- \frac12\mat E)(\mat A - 2\mat E) = \mat E$，
      所以 $\mat A - 2\mat E$ 可逆，
      $(\mat A - 2\mat E) = \frac14(\mat B - 2\mat E)$.
    \end{proof}
    \item $
      \mat A = \ab[\frac14(\mat B - 2\mat E)]^{-1} + 2\mat E =
    \begin{bmatrix}
      6 & 4 & 0\\
      8 & 6 & 0\\
      0 & 0 & 6
    \end{bmatrix}$.
  \end{enumerate}
\end{solution}